\documentclass[11pt]{article}

\usepackage[utf8]{inputenc}
\usepackage[T1]{fontenc}
\usepackage[english]{babel}
\usepackage{lmodern}
\usepackage{microtype}

\usepackage{amsmath,amssymb,amsfonts,amsthm,mathtools}
\usepackage{geometry}
\geometry{a4paper, margin=1in}

\usepackage{array,booktabs,longtable}
\usepackage{enumitem}
\usepackage{hyperref}
\hypersetup{
    colorlinks=true,
    linkcolor=blue,
    citecolor=blue,
    urlcolor=blue
}

\newtheorem{theorem}{Theorem}[section]
\newtheorem{lemma}[theorem]{Lemma}
\newtheorem{proposition}[theorem]{Proposition}
\newtheorem{corollary}[theorem]{Corollary}
\newtheorem{remark}[theorem]{Remark}
\newtheorem{definition}[theorem]{Definition}
\newtheorem{conjecture}[theorem]{Conjecture}

\newcommand{\LL}{\operatorname{lcm}}
\newcommand{\NN}{\mathbb N}
\newcommand{\ZZ}{\mathbb Z}
\newcommand{\QQ}{\mathbb Q}
\newcommand{\PP}{\mathbb P}
\newcommand{\EE}{\mathcal E}
\newcommand{\ord}{\operatorname{ord}}
\newcommand{\vp}{v_p}
\newcommand{\Om}{\omega}
\newcommand{\PhiC}{\Phi}

\title{\textbf{A Diophantine and Structural Study of the Erd\H{o}s--Selfridge Prime-Power Variant of Problem \#889}}
\author{Victor Manuel Mesa Solano\\
\small Independent Researcher}
\date{\today}

\begin{document}

\maketitle

\begin{abstract}
We study the prime-power variant introduced by Erd\H{o}s and Selfridge in their
1967 paper on prime factors of consecutive integers.  For
\[
V(n,k)=\#\{p:\;p^{\nu_p(n+k)}>k\},
\qquad
V_1(n)=\max_{k\geq 1}V(n,k),
\]
Erd\H{o}s and Selfridge asked whether the set
\[
\{n\geq 1:\;V_1(n)=1\}
\]
is finite; they identified
\[
1,2,3,4,6,7,8,12,15,16,24,30,48,80
\]
as the known values below \(2500\), and remarked that \(80\) was probably the
largest solution.  The question remains open.

The purpose of this paper is not to claim a resolution.  Instead, we develop
a detailed Diophantine reduction of the problem and record the advances obtained
so far.  We first reformulate the condition in terms of the ordered
prime-power components of \(n+k\), and derive the decomposition
\[
n+k=d_kq_k,\qquad d_k\mid \operatorname{lcm}(1,\ldots,k),
\]
together with the decisive compatibility condition
\[
\gcd(d_i,d_j)\mid |i-j|.
\]
We then give a complete proof that the branch \(d_2=1\) produces only
the seven small values
\[
1,2,3,6,7,15,30,
\]
using Catalan--Mih\H ailescu, Bang--Zsigmondy, and elementary congruences.
For the branch \(d_2=2,d_3=1\), we prove that if the large prime-power component
of \(n+3\) has exponent \(>1\), then the branch collapses to the single
exception \(n=24\).  We then analyze the complementary case in which \(n+3\)
is prime.  If
\[
n+1=P^\alpha,
\qquad
n+2=2Q^\beta,
\]
the exponents are sharply constrained: when \(\alpha>1\), only the two
known values \(n=8,80\) remain, while when \(\alpha=1\) and \(\beta>1\) the
problem reduces to the residual simultaneous-primality family
\[
2\cdot3^{2^s}-1,\qquad
\frac{3^{2^s}+1}{2},\qquad
2\cdot3^{2^s}+1.
\]
We deliberately leave that final family open rather than claiming a
resolution not supplied by the present arguments.

The remaining leaves are also reduced explicitly.  In particular, for
\(d_2=2,d_3=3\) we obtain \(12\mid n\), the dichotomy
\(d_4\in\{1,4\}\), the collapse \(d_4=1\Rightarrow n=12\), and finally
\(d_5\in\{1,5\}\).  The subbranch \(d_5=1\) yields a system of four
consecutive prime-power equations, while \(d_5=5\) yields a chain of
linear relations between prime powers and, at \(k=6\), an explicit exponential
family when \(d_6=2\).

An exact finite computation up to \(5\cdot 10^7\), recorded together with the
research code, finds no new exceptions and exhibits a very rigid hierarchy of
minimal witnesses.  We emphasize throughout the distinction between
proved statements, named-theorem consequences, and computational observations.

\end{abstract}

\section{Research status and scope}

The original Erd\H{o}s--Selfridge Problem \#889 asks whether
\[
v_0(n):=\max_{k\geq 0} v(n,k)\longrightarrow\infty
\qquad (n\to\infty),
\]
where \(v(n,k)\) counts distinct prime factors of \(n+k\) exceeding \(k\).
The problem is currently listed as open, and the formalized statement in the
Google DeepMind \emph{Formal Conjectures} project is likewise marked
\texttt{research open}.  The prime-power variant discussed in this paper is
also explicitly represented there as an open statement.  See
\cite{erdos1967,erdosproblems,formal889}.

Accordingly, the central objective of this manuscript is to build a rigorous
partial theory rather than to advertise a complete solution.  In particular,
all branches that are not actually proved to be impossible are kept visible
in the final decision tree.

There are three different levels of information used below.

\begin{enumerate}[label=(\roman*)]
    \item \textbf{Elementary or named-theorem proofs.}
    These are lemmas proved in the paper, sometimes using standard results such
    as Catalan--Mih\H ailescu or Bang--Zsigmondy.
    \item \textbf{Exact computational observations.}
    These concern only a finite interval and are not used as substitutes for
    an infinite proof.
    \item \textbf{Research targets.}
    These are surviving Diophantine systems for which we do not yet have a
    complete elimination.
\end{enumerate}

This distinction is essential: the problem remains open even though a large
finite range is completely understood computationally.

\section{The Erd\H{o}s--Selfridge prime-power variant}

\subsection{Prime-power components}

For an integer \(m\geq2\), write
\[
m=\prod_{p\mid m}p^{\nu_p(m)}.
\]
For each prime divisor \(p\mid m\), define its full prime-power component by
\[
Q_p(m):=p^{\nu_p(m)}.
\]
Order the distinct components decreasingly:
\[
Q_1(m)>Q_2(m)>\cdots>Q_{\omega(m)}(m).
\]
If \(m\) has only one prime divisor, put \(Q_2(m)=0\).

We study
\[
V(n,k)=\#\{p:\;Q_p(n+k)>k\}
       =\#\{p:\;p^{\nu_p(n+k)}>k\},
\]
and
\[
V_1(n)=\max_{k\geq1}V(n,k).
\]

Thus \(V(n,k)\leq1\) means that among the full prime-power components of
\(n+k\), at most one exceeds \(k\).

\begin{lemma}\label{lem:q2}
For every \(n\geq1\) and \(k\geq1\),
\[
V(n,k)\geq2
\iff
Q_2(n+k)>k.
\]
\end{lemma}

\begin{proof}
By definition, \(V(n,k)\geq2\) exactly when at least two distinct
prime-power components of \(n+k\) exceed \(k\).  This is equivalent to the
second largest such component exceeding \(k\).
\end{proof}

\subsection{Interval formulation}

Writing \(m=n+k\), the inequality \(Q_2(m)>k\) becomes
\[
m-n>0,\qquad m-n<Q_2(m).
\]
Equivalently,
\[
m-Q_2(m)+1\leq n\leq m-1.
\]
Hence every integer \(m\) contributes an interval
\[
\boxed{
m-Q_2(m)+1\leq n\leq m-1
}
\tag{1}\label{eq:interval}
\]
on which \(V(n,m-n)\geq2\).

Consequently, proving that \(V_1(n)\geq2\) for all sufficiently large \(n\)
is equivalent to proving that these intervals cover all sufficiently large
integers.

\subsection{A universal bound on the offset}

\begin{lemma}\label{lem:kbound}
If \(V(n,k)\geq2\), then
\[
k^2-k<n.
\]
\end{lemma}

\begin{proof}
If \(V(n,k)\geq2\), there are two distinct prime-power components
\(A,B>k\) dividing \(n+k\).  Since they correspond to distinct primes,
they are coprime, so
\[
AB\mid n+k.
\]
Therefore
\[
n+k\geq AB>k^2,
\]
which gives
\[
n>k^2-k.
\]
\end{proof}

Thus for a finite computation \(n\leq N\), it suffices to consider
\[
1\leq k\leq K(N),
\qquad
K(N)=\max\{k:\;k^2-k<N\}.
\]

\section{Exact finite computation}

The accompanying exact sieve computes the second-largest prime-power
component \(Q_2(m)\) for every
\[
m\leq N+K(N)
\]
using a smallest-prime-factor table.  The interval representation
\eqref{eq:interval} is then accumulated with a difference array.

For
\[
N=50\,000\,000,
\qquad
K(N)=7071,
\]
the reported exact computation finds precisely
\[
\boxed{
\EE=
\{1,2,3,4,6,7,8,12,15,16,24,30,48,80\}
}
\tag{2}
\]
as the values with \(V_1(n)=1\).

This is a \emph{finite computational theorem}: it asserts the complete answer
inside the stated range and nothing about larger \(n\).

The same computation defines the minimal witness
\[
\kappa(n):=\min\{k\geq1:\;V(n,k)\geq2\}
\]
for every non-exceptional \(n\leq50\,000\,000\).  The recorded distribution is
\[
\begin{array}{c|r}
\kappa(n)&\text{number of }n\\
\hline
1&46\,997\,805\\
2&2\,876\,654\\
3&113\,669\\
4&11\,288\\
5&484\\
6&80\\
7&4\\
8&2
\end{array}
\tag{3}\label{eq:decomp}
\]
and the last six integers surviving the tests \(k\leq6\) are
\[
240,\quad
471240,\quad
5516280,\quad
16831080,\quad
18164160,\quad
29743560.
\tag{4}
\]
Four are eliminated at \(k=7\), while
\[
\kappa(5516280)=\kappa(18164160)=8.
\]

An important warning is worth recording.  The number \(240\) is \emph{not}
an exception to the prime-power variant: indeed
\[
240+7=247=13\cdot19,
\]
so \(V(240,7)=2\).  It is only an exceptionally late finite-range survivor
of the tests \(k\leq6\).

\section{The divisor decomposition}

From now on, assume
\[
V_1(n)=1.
\]
For each \(k\), put
\[
L_k:=\operatorname{lcm}(1,2,\ldots,k).
\]

All but at most one prime-power component of \(n+k\) are \(\leq k\).  Their
product therefore divides \(L_k\).  Hence we may write
\[
\boxed{
n+k=d_kq_k,
\qquad
d_k\mid L_k,
\qquad
\gcd(d_k,q_k)=1,
}
\tag{5}
\]
where \(q_k\) is either \(1\) or the unique prime-power component which may
exceed \(k\).

For large \(n\), once \(k\geq8\), \(q_k=1\) is impossible because
\(n+k>L_k\) for the relevant range.

\begin{lemma}\label{lem:gcd}
For \(i\neq j\),
\[
\boxed{
\gcd(d_i,d_j)\mid |i-j|.
}
\tag{6}\label{eq:gcd}
\]
In particular,
\[
\boxed{
\gcd(d_k,d_{k+1})=1.
}
\tag{7}
\]
\end{lemma}

\begin{proof}
Suppose a prime \(\ell\) divides both \(d_i\) and \(d_j\).  Then
\(\ell\mid n+i\) and \(\ell\mid n+j\), hence
\[
\ell\mid (n+i)-(n+j)=i-j.
\]
The same statement prime-by-prime gives the divisibility claim for the gcd.
The consecutive case follows because \(|k-(k+1)|=1\).
\end{proof}

The simple relation \eqref{eq:gcd} is the basic mechanism behind the
Diophantine decision tree.

\section{Complete treatment of the branch \(d_2=1\)}

We now prove a substantial theorem which closes the entire \(d_2=1\)
branch.

\subsection{Two consecutive prime powers}

Since \(d_1=1\), we have
\[
n+1=p^a.
\]
If \(d_2=1\), then
\[
n+2=q^b.
\]
Thus \(n+1\) and \(n+2\) are consecutive prime powers.

One of the two numbers is even, and therefore one of the underlying primes
is \(2\).

\subsubsection{Case I: \(n+1=2^a\)}

Then
\[
n+2=2^a+1=q^b.
\]

If \(a,b>1\), Catalan--Mih\H ailescu gives the unique pair of consecutive
perfect powers \(8,9\), hence
\[
n=7.
\]
If \(a=1\), then \(n=1\).

Suppose \(b=1\).  Then
\[
2^a+1
\]
is prime.  For \(a>1\), \(a\) must be a power of \(2\): if
\(a=rs\) with \(r>1\) odd, then
\[
2^a+1=(2^s)^r+1
\]
is composite.  Hence
\[
a=2^t,\qquad t\geq1.
\]
The values \(t=1,2\) give
\[
n=3,\qquad n=15.
\]
It remains to eliminate
\[
\boxed{
n=2^{2^t}-1,\qquad t\geq3.
}
\tag{8}\label{eq:doubleexp}
\]

\begin{lemma}\label{lem:doubleexp}
For every \(t\geq3\),
\[
n=2^{2^t}-1
\]
satisfies \(V(n,k)\geq2\) for some \(k\leq4\).
\end{lemma}

\begin{proof}
Put
\[
M=2^t.
\]

Suppose first that \(t\) is even.  Then \(t\geq4\), and
\[
A:=M-1=2^t-1
\]
is an odd composite integer.

Now
\[
n+3=2^M+2=2(2^{A}+1).
\]
Choose a proper divisor \(u\mid A\) with \(u\geq5\).  Such a divisor exists
because \(A\geq15\) is odd composite.  Then
\[
2^u+1\mid 2^A+1.
\]

By the Bang--Zsigmondy theorem there exists a primitive prime divisor
\(r\) of
\[
2^{2A}-1,
\]
so
\[
\ord_r(2)=2A.
\]
Likewise, since \(u\geq5\), there exists a primitive prime divisor
\(\ell\) of
\[
2^{2u}-1,
\]
with
\[
\ord_\ell(2)=2u.
\]
Both \(r\) and \(\ell\) divide \(2^A+1\), while their orders are different,
so \(r\neq\ell\).  Moreover both primes exceed \(3\).  Thus
\(n+3\) has at least two prime-power components \(>3\), and
\[
V(n,3)\geq2.
\]

Now suppose \(t\) is odd.  Then
\[
M=2^t\equiv2\pmod3.
\]
Since \(2\) has order \(3\) modulo \(7\),
\[
2^M+3\equiv2^2+3\equiv4+3\equiv0\pmod7.
\]
Hence
\[
7\mid n+4.
\]
Also \(n+4\) is odd and
\[
n+4=2^M+3\equiv1\pmod3,
\]
so the quotient \((n+4)/7\) is an integer \(>1\) with no prime factor
\(2\) or \(3\).  Hence it has a prime divisor \(>3\), distinct from \(7\).
Consequently
\[
V(n,4)\geq2.
\]
\end{proof}

Thus the family \eqref{eq:doubleexp} contains no counterexample.

\subsubsection{Case II: \(n+2=2^b\)}

Now
\[
n+1=p^a,\qquad n+2=2^b.
\]
If \(a,b>1\), Catalan--Mih\H ailescu would force the pair
\(8,9\), but the larger member here would have to be \(2^b\), impossible.
Hence
\[
a=1.
\]
Therefore
\[
n=2^b-2,
\qquad
p=2^b-1.
\]
Since \(p\) is prime,
\[
\boxed{b\text{ is prime}.}
\tag{9}
\]
The small primes give
\[
b=2,3,5
\quad\Longrightarrow\quad
n=2,6,30.
\]

Suppose \(b>3\).

If
\[
b\equiv3\pmod4,
\]
then
\[
n+4=2^b+2=2(2^{b-1}+1).
\]
Since \(b-1\equiv2\pmod4\),
\[
5\mid 2^{b-1}+1,
\]
while \(3\nmid 2^{b-1}+1\).  Since \(V(n,4)\leq1\), every prime factor of
the odd number \(2^{b-1}+1\) other than possibly \(5\) would be another
prime \(>4\).  Hence
\[
2^{b-1}+1=5^c.
\]
Modulo \(8\), this gives \(c\) even.  Thus
\[
5^c-2^{b-1}=1
\]
is a pair of consecutive perfect powers with both exponents \(>1\), which
contradicts Catalan--Mih\H ailescu.  Therefore
\[
\boxed{b\equiv1\pmod4.}
\tag{10}\label{eq:d3mod4}
\]

We sharpen this.

\begin{lemma}\label{lem:fermatform}
Under \(d_2=1\), if \(n=2^b-2\) with \(b>5\), then
\[
b=2^{2^s}+1
\]
for some \(s\geq2\).
\end{lemma}

\begin{proof}
We have
\[
n+4=2(2^{b-1}+1).
\]
Put
\[
A=b-1.
\]
By \eqref{eq:d3mod4}, \(4\mid A\).

Suppose \(A\) has an odd divisor \(r>1\).  Write
\[
A=rs.
\]
Then \(s\) is also divisible by \(4\), and
\[
2^s+1\mid2^A+1.
\]

A primitive divisor \(\ell\) of
\[
2^{2s}-1
\]
has
\[
\ord_\ell(2)=2s.
\]
A primitive divisor \(R\) of
\[
2^{2A}-1
\]
has
\[
\ord_R(2)=2A.
\]
Both divide \(2^A+1\), and the orders are different, so
\(\ell\neq R\).  Since \(s\geq4\), both primes exceed \(4\), contradicting
\(V(n,4)\leq1\).

Hence \(A\) has no odd divisor and therefore
\[
A=2^t.
\]
Thus
\[
b=2^t+1.
\]
Since \(b\) is prime, \(t\) itself must be a power of \(2\).  Therefore
\[
b=2^{2^s}+1.
\]
\end{proof}

It remains to eliminate these Fermat-prime-shaped values.

\begin{lemma}\label{lem:fermatbranch}
Let
\[
b=2^{2^s}+1,\qquad s\geq2,
\]
and
\[
n=2^b-2.
\]
Then
\[
V(n,6)\geq2.
\]
\end{lemma}

\begin{proof}
Put
\[
M=b-2=2^{2^s}-1.
\]
Since \(s\geq2\), the exponent \(2^s\) is divisible by \(4\), so
\[
5\mid M.
\]
Also \(M\) is odd and \(M>5\).

We have
\[
n+6
=
2^b+4
=
4(2^{b-2}+1)
=
4(2^M+1).
\]

The divisor \(5\mid M\) implies
\[
2^5+1=33\mid2^M+1.
\]
In particular,
\[
11\mid2^M+1,
\]
and \(11>6\).

On the other hand, Bang--Zsigmondy gives a primitive prime divisor \(R\)
of
\[
2^{2M}-1.
\]
Hence
\[
\ord_R(2)=2M,
\]
so \(R>2M>6\), and \(R\mid2^M+1\).

Finally \(R\neq11\), since their multiplicative orders are different:
\[
2M\neq10.
\]
Thus \(n+6\) has two distinct prime-power components exceeding \(6\),
namely \(11\) and \(R\). Hence
\[
V(n,6)\geq2.
\]
\end{proof}

We have therefore obtained the following complete statement.

\begin{theorem}\label{thm:d2one}
If
\[
V(n,k)\leq1
\qquad\text{for all }1\leq k\leq6
\]
and \(d_2=1\), then
\[
\boxed{
n\in\{1,2,3,6,7,15,30\}.
}
\]
In particular, every counterexample with \(n>30\) satisfies
\[
\boxed{d_2=2.}
\]
\end{theorem}

\begin{remark}
This theorem is stronger than the informal ``Catalan + finite check''
argument in earlier drafts.  The infinite families which appear naturally
after applying Catalan are eliminated using the order-theoretic form of the
Bang--Zsigmondy theorem.
\end{remark}

\section{Consequences of \(d_2=2\)}

Assume now
\[
d_2=2.
\]
Then
\[
n+2=2q_2
\]
with \(q_2\) odd. Hence
\[
n=2(q_2-1),
\]
and therefore
\[
\boxed{4\mid n.}
\tag{11}
\]

Since \(d_3\mid6\) and \(\gcd(d_2,d_3)=1\),
\[
\boxed{d_3\in\{1,3\}.}
\tag{12}
\]

We now treat these two branches separately.

\section{The branch \(d_2=2,\ d_3=1\)}

Write
\[
n+3=r^c.
\]
Because \(4\mid n\),
\[
r^c\equiv3\pmod4,
\]
so
\[
\boxed{
r\equiv3\pmod4,
\qquad
c\text{ odd}.
}
\tag{13}
\]

\subsection{The \(k=4\) reduction}

Set
\[
s=\vp(r+1).
\]
Since \(c\) is odd, the \(2\)-adic form of LTE gives
\[
\vp(r^c+1)=\vp(r+1)=s.
\tag{14}
\]

\subsubsection{Case \(s=2\)}

Then
\[
r^c+1=4A,
\qquad
A\text{ odd}.
\]
Because \(V(n,4)\leq1\), the odd part \(A\) can have at most one prime divisor
exceeding \(4\).  Since \(A\) itself is odd and has no factors \(2\), the
large-prime condition forces \(A\) to be a power of one odd prime whenever
\(A>1\).

Factor
\[
A=
\frac{r+1}{4}
\cdot
\frac{r^c+1}{r+1}
=
A_1B,
\]
where
\[
A_1=\frac{r+1}{4},
\qquad
B=r^{c-1}-r^{c-2}+\cdots-r+1.
\]

If \(r>3\) and \(c>1\), then \(A_1>1\) and \(B>1\).  Since
\(A=A_1B\) is a power of one prime \(\ell\),
\[
A_1=\ell^u,\qquad B=\ell^v.
\]
In particular \(\ell\mid r+1\) and \(\ell\mid B\).  Since
\(r\equiv-1\pmod{r+1}\),
\[
B\equiv c\pmod{r+1},
\]
hence
\[
\ell\mid c.
\]
By LTE,
\[
v_\ell(B)=v_\ell(c).
\]
But \(B=\ell^{v_\ell(c)}\), and therefore
\[
B\leq c.
\]
On the other hand, for \(r\geq3\) and \(c\geq3\),
\[
B=r^{c-1}-r^{c-2}+\cdots-r+1>c.
\]
This contradiction shows:

\[
\boxed{
s=2,\ c>1
\Longrightarrow
r=3.
}
\tag{15}\label{eq:r3reduction}
\]

\subsubsection{Case \(s\geq3\)}

Now \(2^s>4\) itself is a large prime-power component of \(n+4\).  Therefore
the remaining odd part can involve no prime \(>4\).  Since it is odd, the
only possibility is \(1\) or \(3\):
\[
r^c+1=2^s
\qquad\text{or}\qquad
r^c+1=3\cdot2^s.
\]

If
\[
r^c+1=2^s
\]
with \(c>1\), then
\[
r^c+1=(r+1)
(r^{c-1}-r^{c-2}+\cdots-r+1)
\]
has an odd factor \(>1\), impossible.

If
\[
r^c+1=3\cdot2^s,
\]
then
\[
(r+1)
(r^{c-1}-r^{c-2}+\cdots-r+1)=3\cdot2^s.
\]
The second factor is odd and \(>1\), so it must equal \(3\).  But for
\(r\geq3\) and \(c\geq3\),
\[
r^{c-1}-r^{c-2}+\cdots-r+1
\geq r^2-r+1\geq7,
\]
again impossible.

Therefore
\[
\boxed{
s\geq3,\ c>1
\Longrightarrow
\text{contradiction}.
}
\tag{16}\label{eq:slarge}
\]

Combining \eqref{eq:r3reduction} and \eqref{eq:slarge} gives the key reduction.

\begin{lemma}\label{lem:r3}
If
\[
d_2=2,\qquad d_3=1,\qquad c>1,
\]
then
\[
\boxed{r=3}.
\]
Consequently
\[
\boxed{n=3^c-3.}
\]
\end{lemma}

\subsection{The exponent \(c\) is prime}

Now
\[
n+2=3^c-1=2q_2,
\]
so
\[
Q_c:=\frac{3^c-1}{2}
\]
is a prime power.

\begin{lemma}\label{lem:cprime}
If \(c>1\), then \(c\) is prime.
\end{lemma}

\begin{proof}
Since \(c\) is odd, suppose
\[
c=ab,\qquad a,b>1.
\]
Then
\[
Q_c=
\frac{3^a-1}{2}
(1+3^a+\cdots+3^{a(b-1)})
=:AB.
\]
Both \(A\) and \(B\) exceed \(1\).  Since \(AB\) is a prime power,
\[
A=\ell^u,\qquad B=\ell^v
\]
for one prime \(\ell\).

As \(\ell\mid3^a-1\),
\[
B\equiv b\pmod{3^a-1},
\]
so
\[
\ell\mid b.
\]
By LTE,
\[
v_\ell(B)=v_\ell(b).
\]
Thus
\[
B=\ell^{v_\ell(b)}\leq b.
\]
But
\[
B=1+3^a+\cdots+3^{a(b-1)}>b,
\]
contradiction.
\end{proof}

Hence
\[
\boxed{c=p}
\]
for some odd prime \(p\), and
\[
\boxed{n=3^p-3.}
\tag{17}
\]

At this point,
\[
n+1=3^p-2=P^\alpha,
\tag{18}
\]
\[
n+2=3^p-1=2Q^\beta,
\tag{19}
\]
and
\[
n+4=3^p+1=4R^\gamma.
\tag{20}
\]
In particular
\[
\frac{3^p-1}{2}=\PhiC_p(3),
\qquad
\frac{3^p+1}{4}=\PhiC_{2p}(3),
\]
so the \(k=2\) and \(k=4\) conditions impose prime-power values on two
cyclotomic evaluations.

\subsection{The decisive split modulo \(4\)}

\subsubsection{Case \(p\equiv1\pmod4\)}

Consider
\[
n+6=3^p+3=3(3^{p-1}+1).
\]
Because \(p-1\) is even,
\[
\vp(3^{p-1}+1)=1.
\]
Also
\[
3^{p-1}+1\equiv2\pmod3
\]
and, since \(p-1\equiv0\pmod4\),
\[
3^{p-1}+1\equiv2\pmod5.
\]
Therefore
\[
T:=\frac{3^{p-1}+1}{2}
\]
is odd and has no prime factors \(\leq5\).

Since \(V(n,6)\leq1\),
\[
\boxed{T\text{ is a prime power}.}
\tag{21}\label{eq:Tk6}
\]

Write
\[
p-1=2^tm,
\qquad
m\text{ odd}.
\]
Then
\[
T=\frac{3^{2^tm}+1}{2}.
\]
The factor
\[
A=\frac{3^{2^t}+1}{2}
\]
divides \(T\).

Let \(\ell\) be a prime divisor of \(A\).  Then
\[
\ord_\ell(3)\mid 2^{t+1}.
\]
By Bang--Zsigmondy, applied to
\[
3^{2(p-1)}-1,
\]
there exists a primitive prime divisor \(R\) with
\[
\ord_R(3)=2(p-1)=2^{t+1}m.
\]
It divides
\[
3^{p-1}+1,
\]
hence \(T\).  If \(m>1\), then
\[
2^{t+1}m>2^{t+1},
\]
so \(R\neq\ell\).  Thus \(T\) has at least two distinct prime divisors,
contradicting \eqref{eq:Tk6}.  Hence
\[
m=1,
\]
so
\[
p-1=2^t.
\tag{22}
\]

Since \(p>3\) is prime, \(t\) is even.  Thus
\[
p=2^t+1\equiv2\pmod3.
\]
Modulo \(13\), since \(3^3\equiv1\),
\[
3^p+4\equiv3^2+4\equiv0\pmod{13}.
\]
Hence
\[
13\mid n+7.
\]
Moreover \(n+7\) is not divisible by \(2,3,5,7\):
\[
n+7=3^p+4,
\]
\[
3^p+4\equiv1\pmod3,
\]
\[
3^p+4\equiv2\pmod5,
\]
and \(p\equiv5\pmod6\) gives
\[
3^p+4\equiv3^5+4\equiv2\pmod7.
\]
Therefore, under \(V(n,7)\leq1\), the component \(13\) is the only possible
large component, so
\[
3^p+4=13^\delta.
\]
But modulo \(8\),
\[
3^p+4\equiv3+4\equiv7\pmod8,
\]
whereas
\[
13^\delta\equiv1\text{ or }5\pmod8.
\]
Contradiction.

Thus
\[
\boxed{
p\equiv1\pmod4
\quad\Longrightarrow\quad
\text{no counterexample}.
}
\tag{23}
\]

\subsubsection{Case \(p\equiv3\pmod4\)}

Since \(p\equiv3\pmod4\),
\[
3^p-2\equiv2-2\equiv0\pmod5.
\]
But \(n+1=3^p-2=P^\alpha\) is a prime power. Therefore
\[
P=5,
\]
so
\[
3^p-2=5^\alpha.
\tag{24}
\]

Modulo \(8\),
\[
3^p-2\equiv3-2\equiv1\pmod8.
\]
Since
\[
5^\alpha\equiv1\pmod8
\]
only when \(\alpha\) is even,
\[
\alpha=2a.
\]
Thus
\[
(5^a)^2+2=3^p.
\]

We now invoke Ljunggren's classical theorem on the generalized
Lebesgue--Nagell equation: for integers \(x>0\), \(y>1\), \(m>1\),
\[
x^2+2=y^m
\]
has only the solution
\[
(x,y,m)=(5,3,3).
\]
See \cite{ljunggren1943}.

Therefore
\[
a=1,\qquad p=3.
\]
Consequently
\[
\boxed{n=3^3-3=24.}
\]

We have thus proved:

\begin{proposition}\label{prop:branchclosed}
If
\[
d_2=2,\qquad d_3=1,
\]
and the exponent in
\[
n+3=r^c
\]
satisfies \(c>1\), then
\[
\boxed{n=24.}
\]
\end{proposition}

This closes the entire \(d_2=2,d_3=1,c>1\) branch.

\section{Analysis and reduction of the branch $d_2=2,\ d_3=1$}

We now strengthen the analysis of the leaf $d_2=2,d_3=1$.  The proper
prime-power case $n+3=r^c$, $c>1$, was settled in
Proposition~\ref{prop:branchclosed}.  We complete the complementary case
$c=1$ here.

Assume throughout this section that
\[
V_1(n)=1,
\qquad d_2=2,
\qquad d_3=1,
\qquad n+3=R
\]
with $R$ prime.  Then
\[
n+1=P^\alpha,
\qquad
n+2=2Q^\beta,
\]
where $P,Q$ are odd primes.  Equivalently,
\[
\boxed{
P^\alpha+1=2Q^\beta,
\qquad
P^\alpha+2=R.
}
\tag{25}\label{eq:c1system}
\]
The condition $d_2=2$ implies $4\mid n$.

\begin{proposition}[Exponent rigidity on the prime $n+3$ leaf]
\label{prop:primeleaf}
Assume
\[
V_1(n)=1,
\qquad d_2=2,
\qquad d_3=1,
\qquad n+3=R
\]
with $R$ prime.  Write
\[
n+1=P^\alpha,
\qquad
n+2=2Q^\beta.
\]
Then the following reductions hold.
\begin{enumerate}[label=(\roman*)]
    \item If $\alpha>1$, then
    \[
    P=3,
    \qquad \beta=1,
    \qquad \alpha\in\{2,4\},
    \]
    and hence
    \[
    n\in\{8,80\}.
    \]
    \item If $\alpha=1$ and $\beta>1$, then
    \[
    Q=3,
    \qquad \beta=2^s\quad(s\geq1),
    \]
    and necessarily
    \[
    2\cdot3^\beta-1,
    \qquad
    \frac{3^\beta+1}{2},
    \qquad
    2\cdot3^\beta+1
    \]
    are all prime.  The case $s=1$ gives $n=16$.
    The elimination of $s>1$ is left as an explicit residual Diophantine
    problem.
    \item If $\alpha=\beta=1$, then $n=4$.
\end{enumerate}
Consequently, apart from the explicit residual family in (ii), the
$d_2=2,d_3=1$ branch consists only of
\[
\boxed{n\in\{4,8,16,24,80\}}.
\]
\end{proposition}

\begin{proof}
Suppose first that an odd prime $s$ divides $\alpha$, and write
$\alpha=st$ with $s>1$ odd.  Then
\[
P^\alpha+1
=(P^t+1)
\left(P^{t(s-1)}-P^{t(s-2)}+\cdots-P^t+1\right).
\]
The second factor is odd and greater than $1$.  Since the product is
$2Q^\beta$, the odd part of the first factor and the second factor can
involve only the prime $Q$.  Put $X=P^t$ and
\[
B=\frac{X^s+1}{X+1}.
\]
Then $Q\mid X+1$ and, reducing modulo $X+1$,
\[
B\equiv s\pmod{X+1}.
\]
Hence $Q\mid s$.  Since $s$ is odd, LTE gives
\[
\nu_Q(B)=\nu_Q(s).
\]
Thus $B\le s$, whereas $B>s$ for $X\ge3$ and $s\ge3$, a contradiction.
Therefore
\[
\boxed{\alpha=2^u}.
\]

Assume now $\alpha>1$ and put
\[
X=P^{\alpha/2}.
\]
Then
\[
X^2+1=2Q^\beta,
\qquad
R=X^2+2.
\]
Since $R$ is prime, $P\neq3$ would give $X^2\equiv1\pmod3$ and hence
$R\equiv0\pmod3$, impossible for $R>3$.  Thus
\[
P=3.
\]
If $\beta>2$, the classification of
\[
x^2+1=2y^m
\]
for $xy>1$ and $m>2$ gives only $(x,y,m)=(239,13,4)$;
see \cite{Cohn1993,SteinerTzanakis1991}.  This is impossible because
$X$ is a power of $3$.  If $\beta=2$, then $3\mid X$ and reduction modulo
$3$ gives
\[
1\equiv2Q^2\pmod3,
\]
which is impossible.  Hence
\[
\boxed{\beta=1}.
\]

Consequently
\[
n=3^\alpha-1,
\qquad \alpha=2^u.
\]
Suppose $\alpha>2$.  Since
\[
n+4=3(3^{\alpha-1}+1),
\]
the condition $V(n,4)\le1$ implies that
\[
T=\frac{3^{\alpha-1}+1}{4}
\]
is a prime power.  Put $m=\alpha-1=2^u-1$.  If $m=rs$ with $r,s>1$ odd,
then, with $x=3^r$,
\[
T=\frac{x+1}{4}
\left(x^{s-1}-x^{s-2}+\cdots-x+1\right).
\]
The two factors are $>1$.  If their product is a power of one odd prime
$\ell$, then $\ell\mid x+1$ and the second factor is congruent to $s$
modulo $x+1$.  Hence $\ell\mid s$, and LTE gives its $\ell$-adic
valuation as $\nu_\ell(s)$.  It is therefore at most $s$, whereas the
factor itself is $>s$.  Contradiction.  Hence $m$ is prime, and so
$u$ is prime.

If $u$ is odd and $u\ge3$, then $\alpha=2^u$ satisfies
$\alpha\equiv0\pmod4$ and $\alpha\equiv2\pmod3$.  Consequently
\[
5\mid 3^\alpha+4,
\qquad
13\mid 3^\alpha+4.
\]
Let
\[
M=3^\alpha+4.
\]
Since $M\equiv7\pmod8$, $M$ cannot have the form $5^a13^b$, because
$5^a13^b\equiv1$ or $5\pmod8$.  Thus, besides the divisor $13$, there is
another prime divisor exceeding $5$.  The two corresponding prime-power
components both exceed the offset $k=5$, so
\[
V(n,5)\ge2,
\]
a contradiction.  Therefore $u=1$ or $u=2$, giving
\[
\alpha=2,4,
\qquad n=8,80.
\]
This proves (i).

Now assume $\alpha=1$ and $\beta>1$.  Since $P$ and $P+2$ are both prime,
$P\not\equiv1\pmod3$, so $P\equiv2\pmod3$.  Hence
$3\mid P+1=2Q^\beta$, forcing
\[
Q=3.
\]
Thus
\[
P=2\cdot3^\beta-1,
\qquad
P+2=2\cdot3^\beta+1.
\]
If $\beta$ is odd, then according as $\beta\equiv1$ or $3\pmod4$, one of
these two primes is divisible by $5$.  Hence $\beta=2r$ is even.

At $k=4$,
\[
n+4=4\frac{3^\beta+1}{2},
\]
so
\[
T=\frac{3^\beta+1}{2}
\]
is a prime power.  Put $X=3^r$.  Then
\[
X^2+1=2Y^e.
\]
The same classification of $x^2+1=2y^m$ rules out $e>2$ except
$(239,13,4)$, impossible for $X=3^r$; and $e=2$ is impossible modulo
$3$.  Hence $e=1$, so $T$ is prime.

If $r$ has an odd divisor $s>1$, then $3^{2r}+1$ factors as
\[
(3^{2r/s}+1)
\left(3^{2r(s-1)/s}-3^{2r(s-2)/s}+\cdots-3^{2r/s}+1\right),
\]
so $T$ is composite.  Therefore $r=2^w$, and hence
\[
\boxed{\beta=2^{w+1}}.
\]
The case $w=0$ gives $\beta=2$, hence $P=17$ and $n=16$.
For $w>0$ the three quantities
\[
2\cdot3^\beta-1,
\qquad
\frac{3^\beta+1}{2},
\qquad
2\cdot3^\beta+1
\]
must all be prime.  We leave this simultaneous-primality family open.
This proves (ii).

Finally suppose $\alpha=\beta=1$.  Put
\[
q=\frac{n+2}{2}.
\]
Then
\[
n+1=2q-1,
\qquad
n+3=2q+1
\]
are both prime.  If $q>3$, one of these two numbers is divisible by $3$
and is strictly larger than $3$.  Hence $q=3$, giving $n=4$.
\end{proof}

\begin{corollary}[Complete $d_2=2,d_3=1$ branch]
\label{cor:d2d3one}
If $V_1(n)=1$ and
\[
d_2=2,
\qquad d_3=1,
\]
then
\[
\boxed{n\in\{4,8,16,24,80\}}.
\]
Indeed, $c>1$ gives $n=24$ by Proposition~\ref{prop:branchclosed}, while
$c=1$ gives $n\in\{4,8,16,80\}$ by Proposition~\ref{prop:primeleaf}.
\end{corollary}

\begin{remark}[Interpretation]
The corollary classifies an entire state of the $d_k$ decision tree.  Thus
any counterexample with $n>80$ must lie in the branch
\[
\boxed{d_2=2,\qquad d_3=3.}
\]
This concentrates the remaining Diophantine work on the $d_3=3$ branch.
\end{remark}

\subsection{The residual simultaneous-primality family}

The unresolved part of Proposition~\ref{prop:primeleaf} has a particularly
compact one-parameter form.  For $s>1$ put
\[
\beta=2^s.
\]
Then any remaining counterexample on the prime $n+3$ leaf must satisfy
\[
\boxed{
2\cdot3^{2^s}-1,
\qquad
\frac{3^{2^s}+1}{2},
\qquad
2\cdot3^{2^s}+1
\text{ are prime}.
}
\tag{41}\label{eq:residual-prime-triple}
\]
The first and third expressions are the two endpoints $n+1$ and $n+3$,
while the middle expression is forced by the $k=4$ condition.  This turns
the unresolved leaf into a simultaneous-primality problem in a single
power-of-two exponent.  The OEIS records the exponents for which
$2\cdot3^m-1$ is prime and separately those for which $2\cdot3^m+1$ is
prime; these sequences show that the remaining question is not a routine
congruence exercise.  We therefore treat \eqref{eq:residual-prime-triple}
as a genuine residual Diophantine target rather than assert an unsupported
elimination.

\section{The branch \(d_2=2,\ d_3=3\)}

We now assume
\[
d_3=3.
\]
Then
\[
n+3=3q_3,
\]
so \(3\mid n\). Together with \(4\mid n\),
\[
\boxed{12\mid n.}
\tag{26}
\]

Since
\[
d_4\mid\operatorname{lcm}(1,2,3,4)=12
\]
and \(\gcd(d_3,d_4)\mid1\),
we have
\[
3\nmid d_4.
\]
Thus
\[
d_4\in\{1,2,4\}.
\]
But \(4\mid n\) implies \(4\mid n+4\), so the complete \(2\)-adic small
component cannot be represented by \(d_4=2\). Hence
\[
\boxed{d_4\in\{1,4\}.}
\tag{27}
\]

\subsection{The subbranch \(d_4=1\)}

If
\[
d_4=1,
\]
then
\[
n+4=2^a.
\]
Since \(3\mid n\),
\[
2^a\equiv1\pmod3,
\]
so \(a\) is even:
\[
a=2b.
\]
Then
\[
n+3=2^{2b}-1,
\]
and because \(d_3=3\),
\[
\frac{2^{2b}-1}{3}
\]
must be a prime power.

But
\[
\frac{2^{2b}-1}{3}
=
\begin{cases}
(2^b-1)\dfrac{2^b+1}{3},&b\text{ odd},\\[1.2ex]
\dfrac{2^b-1}{3}(2^b+1),&b\text{ even}.
\end{cases}
\tag{28}
\]
For \(b>2\), the two factors are integers \(>1\) and are coprime.  Hence
the product cannot be a prime power.

Therefore
\[
\boxed{d_4=1\Longrightarrow n=12.}
\tag{29}
\]

Consequently, for every \(n>12\) in this branch,
\[
\boxed{d_4=4.}
\tag{30}
\]


\begin{lemma}[The $d_4=4$ branch is divisible by $24$]
\label{lem:24div}
Assume $d_2=2$, $d_3=3$, and $d_4=4$.  Then
\[
\boxed{24\mid n}.
\]
Consequently, if $d_5=5$, then $120\mid n$.
\end{lemma}

\begin{proof}
We know $12\mid n$, so write $n=12m$.  Since
\[
q_4=\frac{n+4}{4}=3m+1
\]
is coprime to $d_4=4$, the integer $q_4$ is odd.  Hence $m$ is even,
so $24\mid n$.  If $d_5=5$, then also $5\mid n$, giving $120\mid n$.
\end{proof}

\subsection{The two subbranches for \(d_5\)}

Since \(12\mid n\),
\[
n+5
\]
is odd and not divisible by \(3\).  Moreover
\[
d_5\mid\operatorname{lcm}(1,\ldots,5)=60
\]
and
\[
\gcd(d_4,d_5)=1.
\]
Hence
\[
\boxed{d_5\in\{1,5\}.}
\tag{31}
\]

\subsubsection{Subbranch \(d_5=1\)}

Here
\[
n+5=p
\]
is prime.  Thus
\[
n+1=p-4=P^\alpha,
\]
\[
n+2=p-3=2Q^\beta,
\]
\[
n+3=p-2=3R^\gamma,
\]
\[
n+4=p-1=4S^\delta.
\]
Hence
\[
\boxed{
\begin{aligned}
p-4&=P^\alpha,\\
p-3&=2Q^\beta,\\
p-2&=3R^\gamma,\\
p-1&=4S^\delta,
\end{aligned}
}
\tag{32}\label{eq:fourpowers}
\]
with \(p\) prime.

Since \(24\mid n=p-5\) by Lemma~\ref{lem:24div},
\[
\boxed{p\equiv5\pmod{24}.}
\tag{33}
\]

The only solution currently visible in the finite computation with the
required $d$-pattern is
\[
p=53,
\qquad n=48.
\]
Indeed,
\[
\begin{aligned}
53-4&=49=7^2,\\
53-3&=50=2\cdot5^2,\\
53-2&=51=3\cdot17,\\
53-1&=52=4\cdot13.
\end{aligned}
\]
The superficially similar value $p=29$ corresponds to $n=24$, but there
$n+3=27=3^3$, so it belongs to the already classified $d_3=1$ branch.

At present we do not have a proof that these are the only solutions to
\eqref{eq:fourpowers}.  This is one of the principal remaining Diophantine targets.

\subsubsection{Subbranch \(d_5=5\)}

Here
\[
\boxed{120\mid n}
\]
by Lemma~\ref{lem:24div}.  Write
\[
n=120u.
\]
Then
\[
\begin{aligned}
n+1&=120u+1=q_1,\\
n+2&=2(60u+1)=2q_2,\\
n+3&=3(40u+1)=3q_3,\\
n+4&=4(30u+1)=4q_4,\\
n+5&=5(24u+1)=5q_5.
\end{aligned}
\tag{34a}\label{eq:qblock120}
\]
Each $q_i$ is a prime power, and the small factors displayed explicitly are
coprime to their corresponding $q_i$.

Subtracting consecutive equations gives the exact prime-power chain
\[
\boxed{
\begin{aligned}
2q_2-q_1&=1,\\
3q_3-2q_2&=1,\\
4q_4-3q_3&=1,\\
5q_5-4q_4&=1.
\end{aligned}}
\tag{34}\label{eq:chain}
\]
This is one of the main Diophantine objects remaining in the project.

\subsection{The $k=6$ split}

Because
\[
n+6=120u+6=6(20u+1),
\]
and $d_6\mid120$ with $\gcd(d_5,d_6)=1$, the only possibilities for a
large counterexample are
\[
\boxed{d_6\in\{2,6\}.}
\tag{35}
\]
Indeed, $d_6=1$ leaves the distinct factors $2$ and $3$ in the remainder,
while $d_6=3$ leaves an even nontrivial remainder and hence cannot produce a
single prime power; the cases containing $5$ are excluded by
$\gcd(d_5,d_6)=1$.

\subsubsection{Case $d_6=6$}

Then
\[
q_6=20u+1
\]
is a prime power and
\[
\boxed{6q_6-5q_5=1.}
\tag{36}\label{eq:q6chain}
\]
Together with \eqref{eq:chain}, this yields the six-term chain
\[
\boxed{
120u+1,\;
60u+1,\;
40u+1,\;
30u+1,\;
24u+1,\;
20u+1
}
\tag{37}\label{eq:qpowers}
\]
of prime-power expressions.  This branch remains unresolved.

\subsubsection{Case $d_6=2$}

Now
\[
n+6=6(20u+1).
\]
Since $d_6=2$, the remaining component is
\[
q_6=3(20u+1).
\]
For this to be a prime power, there must be an integer $b\ge1$ such that
\[
20u+1=3^b.
\]
The congruence modulo $20$ gives
\[
3^b\equiv1\pmod{20},
\]
so
\[
\boxed{b\equiv0\pmod4.}
\]
Write
\[
b=4t,
\qquad
u=\frac{3^{4t}-1}{20}.
\tag{38}\label{eq:expfamily}
\]
Set
\[
N=4t+1.
\]
Then the five prime-power quantities forced by $k=1,\ldots,5$ are
\[
\boxed{
\begin{aligned}
q_1&=2\cdot3^N-5,\\
q_2&=3^N-2,\\
q_3&=2\cdot3^{N-1}-1,\\
q_4&=\frac{3^N-1}{2},\\
q_5&=\frac{2\cdot3^N-1}{5}.
\end{aligned}}
\tag{40}\label{eq:d6two-system}
\]
Each of these five quantities must be a prime power.

\begin{lemma}[Prime-index reduction]
\label{lem:Nprime}
In the branch $d_5=5,d_6=2$, the integer $N=4t+1$ is prime.
\end{lemma}

\begin{proof}
Suppose $N=rs$ with $r,s>1$ odd, and put $x=3^r$.  Then
\[
\frac{3^N-1}{2}
=\frac{x-1}{2}(1+x+\cdots+x^{s-1}).
\]
Both factors exceed $1$.  If their product is a power of one prime $\ell$,
then $\ell$ divides both factors.  Hence $\ell\mid x-1$, while the second
factor is congruent to $s$ modulo $x-1$, so $\ell\mid s$.  By LTE,
\[
\nu_\ell(1+x+\cdots+x^{s-1})=\nu_\ell(s).
\]
Thus the second factor is at most $s$, contradicting
\[
1+x+\cdots+x^{s-1}>s.
\]
Hence $N$ is prime.
\end{proof}

Consequently this leaf reduces to the simultaneous prime-power system
\eqref{eq:d6two-system} with
\[
\boxed{N\equiv1\pmod4,\qquad N\text{ prime}.}
\]

We can sharpen the congruence class of $N$ further.

\begin{lemma}[A mod-$12$ restriction]
\label{lem:Nmod12}
In the branch $d_5=5,d_6=2$, one has
\[
\boxed{N\equiv1\pmod{12}.}
\]
\end{lemma}

\begin{proof}
Since $N$ is prime and $N>3$, we have
\[
N\equiv1\text{ or }5\pmod6.
\]
Suppose $N\equiv5\pmod6$.  Then
\[
N-1\equiv4\pmod6,
\]
and hence, modulo $7$,
\[
3^{N-1}\equiv3^4\equiv4\pmod7.
\]
Therefore
\[
q_3=2\cdot3^{N-1}-1\equiv7\equiv0\pmod7.
\]
But $q_3>7$ for $N>3$, while
\[
q_3\equiv-1\equiv2\pmod3.
\]
Thus $q_3$ cannot be a power of $7$, since every positive power of $7$ is
$1$ modulo $3$.  This contradicts the prime-power condition on $q_3$.
Consequently
\[
N\equiv1\pmod6.
\]
Together with $N\equiv1\pmod4$, this gives
\[
N\equiv1\pmod{12}.
\]
\end{proof}

Thus the current sharp form of the $d_5=5,d_6=2$ branch is
\[
\boxed{
N\text{ prime},\qquad N\equiv1\pmod{12},
}
\]
together with the five simultaneous prime-power conditions
\eqref{eq:d6two-system}.
This is the sharpest form of the branch obtained so far.

\subsection{Targeted computation of the residual $d_6=2$ leaf}

The residual family \eqref{eq:d6two-system} is sufficiently structured that
small exact computations are useful for guiding the next proof.  The table
below records the first few admissible prime indices $N=4t+1$ and a first
prime-power condition that fails.

\[
\begin{array}{c|c|l}
 t&N&\text{first obstruction}\\
\hline
1&5&q_1=481=13\cdot37\\
3&13&q_2=1594321=197\cdot8093\\
4&17&q_1=2\cdot3^{17}-5=13\cdot71\cdot461\cdot607\\
7&29&q_1=13^2\cdot139\cdot5843120971
\end{array}
\]

This table is computational evidence only.  It does, however, illustrate
that the five simultaneous prime-power requirements in
\eqref{eq:d6two-system} are significantly stronger than any one of them
individually.

\section{The current decision tree}

The proved reductions can be summarized as follows.

\[
\begin{array}{c}
V_1(n)=1\\[1mm]
\downarrow\\
d_2=1\quad\text{or}\quad d_2=2\\
\end{array}
\]

The first branch is completely closed:

\[
\boxed{
d_2=1
\Longrightarrow
n\in\{1,2,3,6,7,15,30\}.
}
\]

For \(n>30\),
\[
d_2=2.
\]

Then
\[
d_3\in\{1,3\}.
\]

For \(d_3=1\), write
\[
n+3=r^c.
\]
The case \(c>1\) is completely closed:
\[
\boxed{
d_2=2,\ d_3=1,\ c>1
\Longrightarrow n=24.
}
\]

Therefore every large unresolved point on this branch must have
\[
\boxed{n+3\text{ prime}.}
\]

For \(d_3=3\),
\[
12\mid n.
\]
Then
\[
d_4\in\{1,4\},
\]
with
\[
d_4=1\Longrightarrow n=12.
\]
Hence for \(n>12\),
\[
d_4=4,
\]
and then
\[
d_5\in\{1,5\}.
\]

The resulting research tree is therefore

\[
\boxed{
\begin{array}{rcl}
d_2=1
&\Longrightarrow&
n\in\{1,2,3,6,7,15,30\},\\[1mm]
d_2=2,\ d_3=1,\ c>1
&\Longrightarrow&
n=24,\\[1mm]
d_2=2,\ d_3=1,\ c=1
&\Longrightarrow&
n\in\{4,8,16,80\}\text{ or the residual family \eqref{eq:residual-prime-triple}},\\[1mm]
d_2=2,\ d_3=3,\ d_4=1
&\Longrightarrow&
n=12,\\[1mm]
d_2=2,\ d_3=3,\ d_4=4,\ d_5=1
&\Longrightarrow&
\text{system \eqref{eq:fourpowers}},\\[1mm]
d_2=2,\ d_3=3,\ d_4=4,\ d_5=5
&\Longrightarrow&
\text{chain \eqref{eq:chain} and families \eqref{eq:qpowers}--\eqref{eq:expfamily}}.
\end{array}
}
\tag{39}
\]

This is the main structural output of the present work.

\section{Cyclotomic interpretation}

The branch
\[
d_2=2,\qquad d_3=1,\qquad c>1
\]
provides a particularly clear cyclotomic interpretation.

After proving
\[
n=3^p-3,
\]
the conditions at \(k=2\) and \(k=4\) become
\[
\frac{3^p-1}{2}=\PhiC_p(3)
\]
and
\[
\frac{3^p+1}{4}=\PhiC_{2p}(3).
\]
Thus a counterexample would force both
\[
\PhiC_p(3)
\quad\text{and}\quad
\PhiC_{2p}(3)
\]
to be prime powers, while simultaneously
\[
3^p-2
\]
would have to be a prime power.

The proof above shows how the requirement at \(k=6\) forces strong restrictions
on the exponent \(p\), and the requirement at \(k=7\) then creates the
congruence obstruction
\[
13\mid3^p+4.
\]
In the remaining residue class, the divisibility
\[
5\mid3^p-2
\]
converts the problem into the classical equation
\[
x^2+2=3^p.
\]

This is a useful model for future work: large counterexamples may be forced
into simultaneous primality or prime-power conditions on several nearby
cyclotomic values, after which primitive-divisor theory becomes relevant.

\section{Proof ledger: what is proved and what is not}

For clarity, we record the logical status of the main statements.

\subsection*{Proved in this manuscript}

\begin{enumerate}[label=\arabic*.]
    \item The prime-power-component reformulation
    \[
    V(n,k)\geq2\iff Q_2(n+k)>k.
    \]
    \item The offset bound
    \[
    V(n,k)\geq2\Longrightarrow k^2-k<n.
    \]
    \item The divisor decomposition
    \[
    n+k=d_kq_k,\qquad d_k\mid L_k.
    \]
    \item The gcd compatibility
    \[
    \gcd(d_i,d_j)\mid|i-j|.
    \]
    \item The complete elimination of \(d_2=1\), apart from
    \[
    1,2,3,6,7,15,30.
    \]
    \item The implication
    \[
    d_2=2\Longrightarrow4\mid n.
    \]
    \item The implication
    \[
    d_2=2,\ d_3=1,\ c>1
    \Longrightarrow r=3,\ c\text{ prime}.
    \]
    \item The complete elimination of that latter branch down to
    \[
    n=24.
    \]
    \item The implications
    \[
    d_2=2,\ d_3=3\Longrightarrow12\mid n,
    \]
    \[
    d_4\in\{1,4\},
    \]
    and
    \[
    d_4=1\Longrightarrow n=12.
    \]
    \item For \(n>12\) on the \(d_3=3\) branch,
    \[
    d_4=4,\qquad d_5\in\{1,5\}.
    \]
    \item The prime-power reduction for the \(d_5=1\) and \(d_5=5\) leaves.
    \item The exponential family arising from \(d_6=2\).
\end{enumerate}

\subsection*{Computationally established only}

The exact finite computation up to
\[
5\cdot10^7
\]
and the distribution in Table~1 are computational statements.  They provide
evidence and useful targets, but they do not prove finiteness.

\subsection*{Still open within the present project}

The principal unresolved systems are:

\begin{enumerate}[label=\arabic*.]
    \item The four-consecutive-prime-power system
    \[
    p-4=P^\alpha,\quad
    p-3=2Q^\beta,\quad
    p-2=3R^\gamma,\quad
    p-1=4S^\delta.
    \]
    \item The prime-power chain
    \[
    120u+1,\ 60u+1,\ 40u+1,\ 30u+1,\ 24u+1,\ 20u+1
    \]
    together with the relations \eqref{eq:chain} and \eqref{eq:q6chain}.
    \item The exponential family \eqref{eq:expfamily}.
\end{enumerate}

These are the natural next targets for further work.

\section{Research directions}

The present reductions suggest several concrete directions.

\subsection{Classifying the small factors \(d_k\)}

The condition
\[
\gcd(d_i,d_j)\mid|i-j|
\]
can be interpreted as a compatibility constraint on the prime support of the
sequence
\[
d_1,d_2,d_3,\ldots.
\]
For small \(k\), each \(d_k\) is a divisor of a known least common multiple.
The goal is to classify all compatible patterns before introducing the
large prime-power components \(q_k\).

This produces a finite-state type of problem: each prime \(\ell\) can enter
the small part only for a controlled range of offsets, and repeated use of
the same prime is forced to respect the differences \(i-j\).

\subsection{Primitive divisors of the large components}

The successful elimination of several infinite families used the following
principle.

Suppose
\[
A\mid B
\]
are exponents and both
\[
a^A+1,\qquad a^B+1
\]
occur inside the same prime-power constraint.  A primitive divisor at level
\(2B\) has multiplicative order \(2B\), and therefore cannot coincide with a
primitive divisor at a proper divisor level.

This makes Bang--Zsigmondy particularly well suited to the present problem:
prime-power hypotheses are incompatible with the existence of primitive
divisors at two different levels.

\subsection{Prime-power chains}

The relation
\[
a_1q_1-a_2q_2=1
\]
between prime powers is extremely restrictive.  The branch \(d_5=5\) produces
four such relations immediately, and \(k=6\) adds another.  One possible
approach is to combine congruences modulo carefully chosen primes with
valuation arguments and primitive divisor results.

\subsection{The consecutive-prime-power leaf}

The system
\[
p-4=P^\alpha,\quad
p-3=2Q^\beta,\quad
p-2=3R^\gamma,\quad
p-1=4S^\delta
\]
remains one of the cleanest surviving Diophantine problems in the
$d_3=3,d_4=4,d_5=1$ branch.  The finite computation identifies the
configuration $p=53$ (equivalently $n=48$) among the small cases.  The value
$p=29$ corresponds instead to $n=24$, for which $n+3=27=3^3$ and therefore
belongs to the already classified $d_3=1$ branch.

A useful next theorem would be a nontrivial restriction on the exponents
\[
\alpha,\beta,\gamma,\delta,
\]
for example a parity theorem, a forced exponent-$1$ statement, or a bound
under natural coprimality assumptions.

\section{Conclusion}

The prime-power variant of Erd\H{o}s--Selfridge Problem \#889 remains open.
The main contribution of the present draft is a rigorous structural reduction,
not a complete resolution.

The divisor decomposition
\[
n+k=d_kq_k
\]
and the compatibility law
\[
\gcd(d_i,d_j)\mid|i-j|
\]
convert the original problem into a controlled collection of exponential
Diophantine configurations.  The entire \(d_2=1\) branch can be closed using
Catalan--Mih\H ailescu and Bang--Zsigmondy, while the substantial branch
\[
d_2=2,\ d_3=1,\ c>1
\]
reduces completely to the exceptional value \(n=24\), with the crucial final
step supplied by the classical equation
\[
x^2+2=3^p.
\]

The remaining problem is therefore considerably narrower than the original
formulation.  The main open tasks are now explicit systems of equations
between prime powers, together with the compatible small-factor sequence
\((d_k)\).  This is the point at which the investigation can move from
global experimentation to the classification of individual Diophantine
leaves.

\section*{Acknowledgements}

The author thanks the literature surrounding the Erd\H{o}s Problems project,
the Formal Conjectures project, and the classical work of Erd\H{o}s,
Selfridge, Bang, Zsigmondy, Catalan, Mih\H ailescu, and Ljunggren, whose
results form the framework of the present investigation.

\begin{thebibliography}{99}

\bibitem{erdos1967}
P.~Erd\H{o}s and J.~L.~Selfridge,
\emph{Some problems on the prime factors of consecutive integers},
Illinois Journal of Mathematics \textbf{11} (1967), 428--430.
Available at:
\url{https://combinatorica.hu/~p_erdos/1967-21.pdf}.

\bibitem{erdosproblems}
T.~F.~Bloom,
\emph{Erd\H{o}s Problem \#889},
Erd\H{o}s Problems, current online status,
\url{https://www.erdosproblems.com/889}.

\bibitem{formal889}
Google DeepMind,
\emph{Formal Conjectures: ErdosProblems/889.lean},
\url{https://github.com/google-deepmind/formal-conjectures/blob/main/FormalConjectures/ErdosProblems/889.lean}.

\bibitem{Cohn1993}
J.~H.~E. Cohn,
\emph{The Diophantine equation $x^2+C=y^n$},
Acta Arithmetica \textbf{65} (1993), 367--381.

\bibitem{SteinerTzanakis1991}
R. Steiner and N. Tzanakis,
\emph{Simplifying the solution of Ljunggren's equation $X^2+1=2Y^4$},
Journal of Number Theory \textbf{37} (1991), 123--132.

\bibitem{ljunggren1943}
W.~Ljunggren,
\emph{\"Uber die Gleichung \(x^2+2=y^n\)},
Arkiv f\"or Matematik, Astronomi och Fysik \textbf{29A} (1943), no.~13,
1--11.

\bibitem{zsigmondy}
K.~Zsigmondy,
\emph{Zur Theorie der Potenzreste},
Monatshefte f\"ur Mathematik und Physik \textbf{3} (1892), 265--284.

\bibitem{catalan}
P.~Mih\u{a}ilescu,
\emph{Primary cyclotomic units and a proof of Catalan's conjecture},
Journal f\"ur die reine und angewandte Mathematik \textbf{572} (2004),
167--195.

\bibitem{bloomreport}
T.~F.~Bloom,
\emph{Erd\H{o}s \#889 -- working report},
2026.
The report distinguishes rigorous arguments from computational observations
and records the current open status of the problem.

\end{thebibliography}

\appendix

\section{Exact computational sieve}

The principal computational object is the second-largest prime-power
component \(Q_2(m)\).  The following implementation uses a linear smallest
prime factor sieve and then accumulates the intervals \eqref{eq:interval}.

\begin{verbatim}
#include <bits/stdc++.h>
using namespace std;

// Exact search for the Erdős–Selfridge prime-power variant:
//
// V(n,k) = number of p such that p^{nu_p(n+k)} > k.
//
// We seek all n <= N with max_{k>=1} V(n,k) <= 1.
//
// Let Q2(m) be the second-largest prime-power component of m.
// Then V(n,k) >= 2 iff Q2(n+k) > k.
// Hence m=n+k covers n precisely when
//
//     m - Q2(m) + 1 <= n <= m-1.
//
// Also, V(n,k)>=2 implies k^2-k<n.

static int Kbound(long long N) {
    long long k = 1;
    while ((k + 1) * (k + 1) - (k + 1) < N) ++k;
    return (int)k;
}

int main(int argc, char** argv) {
    int N = (argc >= 2 ? atoi(argv[1]) : 1000000);
    if (N < 1) return 1;

    const int K = Kbound(N);
    const int M = N + K;

    cerr << "N=" << N << " K=" << K << " M=" << M << '\n';

    // Smallest prime factor table via a linear sieve.
    vector<uint32_t> spf(M + 1, 0);
    vector<int> primes;
    primes.reserve(M / max(1, (int)log((double)M)));

    for (int i = 2; i <= M; ++i) {
        if (spf[i] == 0) {
            spf[i] = (uint32_t)i;
            primes.push_back(i);
        }

        for (int p : primes) {
            if (p > (int)spf[i] || 1LL * i * p > M) break;
            spf[i * p] = (uint32_t)p;
        }
    }

    // q2[m] = second-largest prime-power component of m.
    vector<uint32_t> q2(M + 1, 0);

    for (int m = 2; m <= M; ++m) {
        int x = m;
        uint64_t best1 = 0, best2 = 0;

        while (x > 1) {
            int p = (int)spf[x];
            uint64_t q = p;

            x /= p;
            while (x > 1 && (int)spf[x] == p) {
                x /= p;
                q *= (uint64_t)p;
            }

            if (q > best1) {
                best2 = best1;
                best1 = q;
            } else if (q > best2) {
                best2 = q;
            }
        }

        q2[m] = (uint32_t)best2;
    }

    // Difference-array interval coverage.
    vector<int32_t> diff(N + 2, 0);

    for (int m = 2; m <= M; ++m) {
        uint32_t q = q2[m];
        if (q == 0) continue;

        int L = max(1, m - (int)q + 1);
        int R = min(N, m - 1);

        if (L <= R) {
            ++diff[L];
            --diff[R + 1];
        }
    }

    vector<int> exceptions;
    int coverage = 0;

    for (int n = 1; n <= N; ++n) {
        coverage += diff[n];
        if (coverage == 0) exceptions.push_back(n);
    }

    cout << "exceptions=" << exceptions.size() << '\n';
    cout << "last="
         << (exceptions.empty() ? -1 : exceptions.back()) << '\n';

    for (size_t i = 0; i < exceptions.size(); ++i) {
        if (i) cout << ' ';
        cout << exceptions[i];
    }
    cout << '\n';

    vector<long long> hist(32, 0);
    int max_k = 0, max_n = -1;

    unordered_set<int> is_exception;
    is_exception.reserve(exceptions.size() * 2 + 1);

    for (int n : exceptions) is_exception.insert(n);

    for (int n = 1; n <= N; ++n) {
        if (is_exception.count(n)) continue;

        int witness = -1;

        for (int k = 1; k <= K; ++k) {
            if (q2[n + k] > (uint32_t)k) {
                witness = k;
                break;
            }
        }

        if (witness < 0) {
            cerr << "ERROR: no witness for n=" << n << '\n';
            return 2;
        }

        if (witness >= (int)hist.size())
            hist.resize(witness + 1);

        ++hist[witness];

        if (witness > max_k) {
            max_k = witness;
            max_n = n;
        }
    }

    cout << "max_min_k=" << max_k
         << " at_n=" << max_n << '\n';

    for (int k = 1; k <= max_k; ++k) {
        if (hist[k])
            cout << k << ':' << hist[k] << ' ';
    }

    cout << '\n';
}
\end{verbatim}

\section{Selected arithmetic identities}

For convenient reference, the main identities used in the reductions are
collected here.

\subsection*{Branch \(d_2=1\)}

\[
n+1=2^a,\qquad n+2=2^a+1,
\]
or
\[
n+1=p,\qquad n+2=2^b,
\]
with \(b\) prime.

The two infinite families arising before the Zsigmondy elimination are
\[
n=2^{2^t}-1
\]
and
\[
n=2^b-2,\qquad b\equiv1\pmod4.
\]
The latter reduces further to the Fermat-prime-shaped family
\[
b=2^{2^s}+1.
\]

\subsection*{Branch \(d_2=2,d_3=1,c>1\)}

\[
n=3^p-3,
\]
\[
\frac{3^p-1}{2}=\PhiC_p(3),
\]
\[
\frac{3^p+1}{4}=\PhiC_{2p}(3),
\]
\[
3^p-2=P^\alpha.
\]

The two decisive congruence regimes are
\[
p\equiv1\pmod4
\]
and
\[
p\equiv3\pmod4.
\]

In the first,
\[
p-1=2^t
\]
is forced by the prime-power condition at \(k=6\), followed by
\[
13\mid3^p+4
\]
at \(k=7\).

In the second,
\[
5\mid3^p-2,
\]
hence
\[
3^p-2=5^{2a},
\]
reducing to
\[
(5^a)^2+2=3^p.
\]

\subsection*{Branch \(d_2=2,d_3=3\)}

\[
12\mid n,
\]
\[
d_4\in\{1,4\},
\]
\[
d_4=1\Longrightarrow n=12,
\]
and for \(n>12\),
\[
d_4=4,
\qquad
d_5\in\{1,5\}.
\]

For \(d_5=5\),
\[
n=120u
\]
and
\[
\begin{aligned}
2q_2-q_1&=1,\\
3q_3-2q_2&=1,\\
4q_4-3q_3&=1,\\
5q_5-4q_4&=1.
\end{aligned}
\]

If \(d_6=2\),
\[
20u+1=3^{4t},
\qquad
u=\frac{3^{4t}-1}{20}.
\]

\section{A checklist for the next proof phase}

The natural next step is not a larger blind search.  It is to attack the
remaining leaves one by one.

\begin{enumerate}[label=\arabic*.]
    \item Study the equation
    \[
    P^\alpha+1=2Q^\beta
    \]
    together with the primality of
    \[
    P^\alpha+2.
    \]
    The known values \(n=8,16,80\) should emerge naturally.

    \item Study the four consecutive relations
    \[
    p-4=P^\alpha,\quad
    p-3=2Q^\beta,\quad
    p-2=3R^\gamma,\quad
    p-1=4S^\delta.
    \]
    The first relevant example in this branch is \(p=53\) (equivalently \(n=48\)).

    \item For \(d_5=5\), exploit the linear prime-power chain and try to force
    incompatible residue classes.

    \item In the \(d_6=2\) exponential branch, apply primitive-divisor
    arguments to the numbers
    \[
    3^{4t+1}-2,\quad
    2\cdot3^{4t}-1,\quad
    \frac{3^{4t+1}-1}{2},
    \quad
    \frac{6\cdot3^{4t}-1}{5}.
    \]

    \item Extend the \(d_k\)-classification to \(k=7,8,\ldots\) only after
    the small-\(k\) leaves have been simplified as far as possible.
\end{enumerate}

\end{document}
